The four given directions can be interpreted as one continuous line segment motion rather than four discrete motions. The goal is to gradually develop a more accurate diagnosis of multimodal integration. Starting from differences in behaviors between reasoning categories, then moving to the structure of what each model retrieves from the video, then into inner dynamics of attention, and eventually concludes in causal verification. 
\paragraph{Expansion to the Twelve Reasoning Categories}
The experiments in this thesis treat spatial, temporal, and causal reasoning as the primary axes of analysis, without asking whether modality contribution itself varies across the full set of reasoning phenomena covered by \textit{MAIA}. A natural extension is to apply the same behavioural and modality-contribution metrics separately within each of the twelve categories.
\paragraph{Joint Analysis of Semantic Availability, Attention, and Behaviour}
This direction develops directly out of the attention analysis introduced in this thesis. Rather than treating semantic availability, attention allocation, and behavioural contribution as three separate questions, a natural next step is to study their joint relationship:
\begin{equation}
\text{Semantic Availability}
\;\longleftrightarrow\;
\text{Attention Allocation}
\;\longleftrightarrow\;
\text{Behavioural Contribution}.
\end{equation}
The question is no longer only \textit {whether the model found a given piece of information, how much it attends to a given modality, or how much accuracy improves when that modality is added}, but \textbf{ how these three quantities relate to one another at the level of individual items}. This connects directly back to the question raised by \textit{Lost in Transcription} \cite{deodati2026lost}: whether a modality that is available is also genuinely exploited, and whether its internal use corresponds to an observable contribution to the final decision, rather than being available, attended, and yet behaviourally inert.
